%%\section{\minerva Experiment}

 \minerva uses the neutrinos from the NuMI beamline at Fermi National Accelerator Laboratory~\cite{Adamson:2015dkw}. 
The neutrino flux for the data presented here is peaked at 3 GeV and contains $95\%$ \numu, with the remainder consisting of \numubar, $\nu_e$, and $\bar\nu_e$ ~\cite{Alt:2006fr}. The neutrino beam is simulated with Geant4 9.2.p03~\cite{Agostinelli:2002hh}, and constrained with thin-target hadron production measurements and an in-situ neutrino electron scattering constraint~\cite{Aliaga:2016oaz}. The analysis uses data collected between March 2010 and April 2012, and corresponds to $3.06\times{10}^{20}$ protons on target (POT).

The MINER$\nu$A detector~\cite{Aliaga:2013uqz} is segmented longitudinally into several regions: nuclear targets, the scintillator tracker, and downstream electromagnetic and hadronic calorimeters. The nuclear target region contains five solid passive targets of carbon(C), iron(Fe), and lead(Pb), separated from each other by 4 or 8 scintillator planes for vertex and particle reconstruction. Targets 1, 2 and 3 contain distinct segments of Fe and Pb planes that are 2.6 cm thick; target 3 also has a C segment which is 7.6 cm thick and target 5 has Fe and Pb segments which are 1.3 cm thick. The analysis restricts to targets with 8 scintillator planes on both sides, except target 5 which has 4 scintillator planes upstream the target. The tracker is made solely of scintillator planes. Strips in adjacent planes are rotated by $60^{\circ}$ from each other, which enables three-dimensional track reconstruction~\cite{Aliaga:2013uqz}. The MINOS Near Detector is two meters downstream of the MINER$\nu$A detector and serves as a magnetized muon spectrometer~\cite{Michael:2008bc}.

The neutrino event generator \genie 2.8.4~\cite{Andreopoulos:2009rq} is used to simulate neutrino interactions in the detector. The CCQE scattering model uses a relativistic Fermi gas model (RFG) and dipole axial form factor with $M_A=0.99$ GeV. Resonant production is modeled using the Rein-Sehgal model~\cite{Rein:1980wg}, deep inelastic scattering (DIS) kinematics is modeled using 2003 Bodek-Yang model~\cite{Bodek:2002ps} and the hadron final states are modeled with Koba-Nielsen-Olsen scaling and PYTHIA~\cite{koba,pythia}. The default GENIE simulation has been augmented to include interactions resulting in two particles and two holes (2p2h), as formulated in the Valencia model \cite{Nieves:2011yp,Gran:2013kda,Schwehr:2016pvn}. The relative strength of this 2p2h prediction has been tuned to \minerva inclusive scattering data~\cite{Rodrigues:2015hik}. The RPA effect from the calculation of~\cite{Nieves:2004wx} is included for quasielastic events. Moreover, the GENIE nonresonant pion production prediction has been modified to agree with deuterium data~\cite{Wilkinson:2014yfa}. To treat FSI, \genie uses an effective model, in which hadronic intranuclear rescattering cross sections increase with the nuclear size according to $A^{2/3}$ scaling~\cite{Dytman:2011zz}. The final state distributions in energy and angle come from 2-body kinematics and phase space formulations.  The model provides a good description of hadron-nucleus data. Comparisons are also made to predictions of NuWro event generator~\cite{Golan:2012wx}, which uses a local Fermi gas model, an intranuclear cascade of hadronic interactions in the FSI model and medium corrections~\cite{Salcedo:1987md}. Coulomb corrections are not included in the simulations~\cite{Coulomb}.

 The interactions and decays of particles produced in the neutrino interactions of the final-state particles that exit the nucleus are simulated by Geant4 9.4.2 ~\cite{Agostinelli:2002hh}. The visible energy scale is calibrated using through-going muons, such that the energy deposit (per plane) distribution is the same for data and simulation.  Measurements made with a smaller version of the \minerva detector in a
hadron test beam~\cite{Aliaga:2013uqz} are used to constrain the
uncertainties associated with the detector responses to both protons and charged pions.  
